% BT967 -- selector rail theorem insert for w33_paper.tex

\subsection{Selector rails, phase slots, and packet ABI (BT962--BT967)}\label{subsec:bt967-selector-rails}

The final E8 selector
\[
  [(3,68),(4,42),(38,65),(90,144)]
\]
turns the rank-eight chain shadow into four canonical hyperbolic rails.  Their
support sums and xor masks are
\[
  (12,12,14,22),\qquad (71,46,91,234).
\]
Every exactly-three-rail face has size \(27\).  Choosing the three faces that
contain the high-support rail gives the canonical \(27+27+27\) split
\[
  (0,1,3),\qquad (0,2,3),\qquad (1,2,3),
\]
with complementary \(27\)-face \((0,1,2)\).

For phase-slot bookkeeping we use
\[
  \mathrm{score}(r)=\mathrm{support}(r)+\operatorname{popcount}(\mathrm{xor}(r)).
\]
This gives scores
\[
  16,16,19,27.
\]
The two light rails remain tied by support and phase score.  The packet ABI
breaks this ordering by ascending xor mask: rail \(1\) precedes rail \(0\)
because \(46<71\).  Hence the selector-backed Holonet prefix assignment is
\[
\begin{array}{cclcl}
0   &\mapsto& \text{mirror}   &\mapsto& \text{rail }1,\ (4,42),\\
10  &\mapsto& \text{schedule} &\mapsto& \text{rail }0,\ (3,68),\\
110 &\mapsto& \text{cache}_A  &\mapsto& \text{rail }2,\ (38,65),\\
111 &\mapsto& \text{cache}_B  &\mapsto& \text{rail }3,\ (90,144).
\end{array}
\]
The prefix code is prefix-free.  What remains open is dynamic preservation:
explicit lane-action maps for mirror, schedule, and the two cache roles must be
encoded before packet-lane preservation can be promoted from ABI convention to
theorem.
